\documentclass{article}
\usepackage{amsmath, amssymb, amsthm}
\usepackage{hyperref}
\usepackage{geometry}
\geometry{margin=1in}

\title{Erd\H{o}s Problem \#187}
\date{Last updated: 2025-08-31}
\author{Source: erdosproblems.com}

\begin{document}
\maketitle

\noindent\textbf{Status:} OPEN \quad \textbf{No prize}

\noindent\textbf{Tags:} additive combinatorics, ramsey theory

\medskip

\noindent\textbf{Problem Statement:}

Find the best function $f(d)$ such that, in any 2-colouring of the integers, at least one colour class contains an arithmetic progression with common difference $d$ of length $f(d)$ for infinitely many $d$.




Originally asked by Cohen. Erd\H{o}s observed that colouring according to whether $\{ \sqrt{2}n\}&#60;1/2$ or not implies $f(d) \ll d$ (using the fact that $\|\sqrt{2}q\| \gg 1/q$ for all $q$, where $\|x\|$ is the distance to the nearest integer). In \cite{Er80} he states that Petruska and Szemer\'{e}di have proved that $f(d) \ll d^{1/2}$, and expected $f(d)\leq d^{o(1)}$. Beck \cite{Be80} has improved this using the probabilistic method, constructing a colouring that shows $f(d)\leq (1+o(1))\log_2 d$. Van der Waerden's theorem implies $f(d)\to \infty$ is necessary.




References

[Be80] Beck, J\'{o}zsef, A remark concerning arithmetic progressions . J. Combin. Theory Ser. A (1980), 376-379.

[Er80] Erd\H{o}s, Paul, A survey of problems in combinatorial number theory . Ann. Discrete Math. (1980), 89-115.


\bigskip
\noindent\small{Source: \url{https://www.erdosproblems.com/187}}
\end{document}
